\chapter{Limitações e conclusões}
\section{Alcance da avaliação}
A investigação concentrou-se em uma área de distribuição e em um horizonte temporal curto. Há dependência entre janelas da mesma UC e heterogeneidade entre circuitos. UCs sem alimentador e instalações com telemetria insuficiente seguem percursos distintos de análise. A falta de associação confiável a transformadores restringiu a localização topológica. O modelo também depende da data de disponibilidade dos alarmes e da cadência efetiva de cada canal.

A inferência OOT utilizou um corte posterior ao treino. Esse desenho demonstrou aplicação de uma referência congelada em data futura ao ajuste, sem estabelecer estabilidade entre estações ou em outras áreas. A concordância entre detectores, a retenção sob diferentes coberturas e a permanência em uma lista são propriedades distintas. Pesos heurísticos e escolha da referência influenciam as três. A divergência entre a cobertura horária em grade fixa de 168 horas e a cobertura da janela analítica na UC 685***20 ilustrou a sensibilidade da métrica à definição do denominador temporal. A replicação independente depende dos artefatos congelados e das tabelas de fatos; a sensibilidade dos pesos e dos limiares heurísticos demanda ensaios próprios. Rótulos independentes, novas datas de teste e comparação entre tecnologias de medição são necessários para avaliar a generalização.

\section{Conclusões e continuidade}
A reconstrução da vigência UC–NIO alterou a interpretação física dos casos. Na UC 908***00, reconhecer a continuidade do equipamento permitiu relacionar a queda de FA ao histórico anterior e examinar a reversão de fluxo. Na UC 291***96, o contraste entre cobertura aparente e integridade do acumulador deslocou a investigação para o denominador temporal. Os casos mostram por que a plataforma preserva simultaneamente a série elétrica, a qualidade de observação e o cadastro da instalação.

A hipótese de diferenciação das listas encontrou suporte na composição dos conjuntos. Na execução corrigida, o filtro topológico distinguiu 4.311 candidatos com contexto de circuito de uma lista OOT geral com 6.037 unidades. A retirada do LOF substituiu 187 membros do Top-2.000; a do PCA residual, 166; e o núcleo de quatro detectores, 254. Os quantis congelados preservaram os escores no ensaio de perturbação do lote, enquanto a escala dinâmica produziu deslocamento máximo de 0,2291. Esses resultados sustentam a rastreabilidade da seleção e delimitam a consequência operacional de uma escolha metodológica: duas políticas podem consumir a mesma capacidade e encaminhar UCs diferentes.

A contribuição do trabalho consiste em tornar reconstituível o percurso entre leitura, janela, hipótese diagnóstica e posição na carteira. A avaliação de perdas comprovadas, rendimento das inspeções, energia recuperada e retorno econômico permanece condicionada a OS, TOI e acompanhamento prospectivo independente. Um piloto poderá vincular as listas congeladas aos desfechos e avaliar, sob a mesma capacidade, a repriorização de casos conhecidos e a descoberta incremental.
